%! AP Statistics — Unit 1, Lesson 1.4 — Guided Notes
\documentclass[10pt]{article}
\usepackage{apstats-article}
\usepackage{apstats-boxes}

%=============================================================================
\begin{document}

\pageheader{Unit 1, Lesson 1.4}{Guided Notes}
\namedateperiod

\vspace{0.12in}

% ── OBJECTIVE ────────────────────────────────────────────────────────────────
\begin{objectivebox}
\small By the end of this lesson, I will be able to\dots\\[4pt]
\writeline\\[1pt]
\writeline\\[1pt]
\writeline
\end{objectivebox}

\vspace{0.1in}

% ── VOCABULARY ───────────────────────────────────────────────────────────────
\begin{vocabbox}
\termblanklong{Bar chart (bar graph)}
\termblanklong{Relative frequency bar chart}
\termblanklong{Pie chart}
\termblanklong{Segmented bar chart}
\termblanklong{Misleading graph}
\end{vocabbox}

\vspace{0.1in}

% ── HOOK ─────────────────────────────────────────────────────────────────────
\begin{hookbox}
\small\textbf{What's Wrong With These Graphs? --- Discuss with a partner.}

\vspace{8pt}
\textbf{Graph A} shows monthly sales for two products. The y-axis starts at 950.
Bar for Product X reaches 970; bar for Product Y reaches 1000.
A reader glances at the graph and says ``Product Y sells twice as much!''

\vspace{6pt}
\begin{enumerate}[leftmargin=1.5em, itemsep=6pt]
  \item Is the reader's conclusion correct? Why or why not?\\[5pt]
        \writeline\\[1pt]\writeline

  \item How would you fix Graph A so it is not misleading?\\[5pt]
        \writeline

  \item \textbf{Graph B} is a 3-D pie chart with the largest slice tilted
        toward the viewer. How might this distort the reader's perception?\\[5pt]
        \writeline\\[1pt]\writeline
\end{enumerate}
\end{hookbox}

\vspace{0.1in}

% ── BAR CHART ────────────────────────────────────────────────────────────────
\begin{notesbox}{Constructing a Bar Chart --- Four Requirements}
\small

\begin{multicols}{2}

\textbf{\textcolor{navy}{A bar chart must have:}}
\begin{enumerate}[leftmargin=1.4em, itemsep=4pt]
  \item A \blank{3cm} (what the graph shows).
  \item A \textit{labeled} \blank{2.5cm} axis listing the categories.
  \item A \textit{labeled} \blank{2.5cm} axis with a scale starting at \blank{1cm}.
  \item \blank{2cm} between bars (categories have no numeric order).
\end{enumerate}

\vspace{4pt}
{\color{navy}\small\textit{Note:} The y-axis may show frequency \textit{or}
relative frequency --- label it clearly!}

\columnbreak

\textbf{\textcolor{navy}{Frequency vs.\ Relative Frequency}}

\vspace{4pt}
\begin{tabular}{@{}p{0.42\linewidth} p{0.42\linewidth}@{}}
\textbf{Frequency bar chart} & \textbf{Rel.\ freq.\ bar chart} \\
y-axis: \blank{1.5cm} & y-axis: \blank{2cm} \\[4pt]
Use when: \blank{1.5cm} & Use when: \blank{1.5cm} \\
\blank{2cm} & comparing \blank{1.5cm} \\
\end{tabular}

\end{multicols}
\end{notesbox}

\vspace{0.1in}

% ── WORKED EXAMPLE — BAR CHART ───────────────────────────────────────────────
\begin{notesbox}{Worked Example --- Pet Ownership Bar Chart ($n = 40$)}
\small

\noindent Data from Lesson 1.3: Dog~30\%, Cat~20\%, Fish~15\%, Bird~10\%, None~25\%.
Construct a \textbf{relative frequency bar chart} in the space below.

\vspace{8pt}
\textbf{Label the axes, add a title, and draw the bars. Leave gaps between bars.}

\vspace{10pt}
\begin{tcolorbox}[colback=white, colframe=slate, boxrule=0.5pt, arc=2mm,
    left=3mm, right=3mm, top=3mm, bottom=3mm, height=2.2in]
\end{tcolorbox}

\vspace{10pt}
\textbf{Write one AP-style sentence interpreting the bar chart:}\\[6pt]
\writeline\\[1pt]\writeline

\end{notesbox}

\newpage

% ── PIE CHART ────────────────────────────────────────────────────────────────
\begin{notesbox}{Constructing a Pie Chart}
\small
\begin{multicols}{2}

\textbf{\textcolor{navy}{How to find each sector's angle:}}\\[4pt]
\[
  \text{Sector angle} = \text{relative frequency} \times \blank{2cm}
\]

\vspace{4pt}
\textbf{\textcolor{navy}{Example (Pet data):}}\\[4pt]
\renewcommand{\arraystretch}{1.6}
\begin{tabular}{@{} >{\bfseries}p{0.22\linewidth}
                     >{\centering\arraybackslash}p{0.18\linewidth}
                     >{\centering\arraybackslash}p{0.24\linewidth} @{}}
\rowcolor{sky}
\textbf{Category} & \textbf{Rel.\ Freq.} & \textbf{Angle (${}^\circ$)} \\
\midrule
Dog  & 0.30 & \blank{1.8cm} \\
Cat  & 0.20 & \blank{1.8cm} \\
Fish & 0.15 & \blank{1.8cm} \\
Bird & 0.10 & \blank{1.8cm} \\
None & 0.25 & \blank{1.8cm} \\
\midrule
\rowcolor{sky}\textbf{Total} & 1.00 & \blank{1.8cm} \\
\end{tabular}

\columnbreak

\textbf{\textcolor{navy}{When to use a pie chart:}}
\begin{itemize}[itemsep=3pt]
  \item Emphasizes \blank{2.5cm} relationships (each slice is a part of the whole).
  \item Works best with \blank{1.5cm} categories (6 or fewer).
  \item \textit{Not} ideal for \blank{2.5cm} categories because slices become too small to read.
\end{itemize}

\vspace{6pt}
\textbf{\textcolor{navy}{Bar vs.\ Pie:}}\\[4pt]
Bar charts are generally preferred because they are easier to \blank{2.5cm} and do not rely on judging \blank{2cm} sizes.

\end{multicols}
\end{notesbox}

\vspace{0.1in}

% ── MISLEADING GRAPHS ────────────────────────────────────────────────────────
\begin{notesbox}{Misleading Graphs --- Three Common Tricks}
\small

\begin{multicols}{2}

\textbf{Trick 1: Non-zero y-axis}\\
The y-axis starts above \blank{1cm}, making small differences look large.\\[4pt]
\textit{Fix:} \blank{3.5cm}

\vspace{8pt}
\textbf{Trick 2: 3-D or perspective effects}\\
Slices or bars closest to the viewer appear \blank{2cm}.\\[4pt]
\textit{Fix:} \blank{3.5cm}

\columnbreak

\textbf{Trick 3: Inconsistent scale or categories}\\
Category widths differ, or the scale jumps unevenly.\\[4pt]
\textit{Fix:} \blank{3.5cm}

\vspace{8pt}
\textbf{\textcolor{navy}{AP-style question:}}\\
``Explain one way the graph is misleading and describe how to fix it.''\\[4pt]
Always identify \textit{what} is wrong and \textit{how} it distorts the reader's conclusion.

\end{multicols}
\end{notesbox}

\vspace{0.1in}

% ── GUIDED PRACTICE ──────────────────────────────────────────────────────────
\begin{practicebox}
\small
\textbf{A survey of 80 adults asks: ``What is your favorite TV genre?''}\\
Results: News~32, Sports~20, Reality~16, Other~12.

\vspace{8pt}
\renewcommand{\arraystretch}{1.8}
\begin{tabularx}{\linewidth}{@{} >{\bfseries}p{0.20\linewidth}
                                    C{0.18\linewidth}
                                    C{0.24\linewidth}
                                    C{0.22\linewidth} @{}}
\rowcolor{sky}
\textbf{Genre} & \textbf{Frequency} & \textbf{Rel.\ Frequency} & \textbf{Angle (${}^\circ$)} \\
\midrule
News    & 32 & \blank{2.5cm} & \blank{2cm} \\
Sports  & 20 & \blank{2.5cm} & \blank{2cm} \\
Reality & 16 & \blank{2.5cm} & \blank{2cm} \\
Other   & 12 & \blank{2.5cm} & \blank{2cm} \\
\midrule
\rowcolor{sky}\textbf{Total} & \blank{2cm} & \blank{2.5cm} & \blank{2cm} \\
\end{tabularx}

\vspace{10pt}
\begin{enumerate}[leftmargin=1.5em, itemsep=6pt]
  \item Sketch a \textbf{relative frequency bar chart} in the box below.
        Include a title, labeled axes, scale from 0, and gaps between bars.

\vspace{6pt}
\begin{tcolorbox}[colback=white, colframe=slate, boxrule=0.5pt, arc=2mm,
    left=2mm, right=2mm, top=2mm, bottom=2mm, height=1.8in]
\end{tcolorbox}

  \item Write one sentence describing the distribution of TV genre preference
        for these 80 adults.\\[6pt]
        \writeline\\[1pt]\writeline

  \item Someone redraws the bar chart with the y-axis starting at 10\%.
        How does this change what a reader sees? Is it misleading?\\[6pt]
        \writeline\\[1pt]\writeline
\end{enumerate}
\end{practicebox}

\vspace{0.1in}

% ── SPIRAL / BIG IDEAS ────────────────────────────────────────────────────────
\begin{spiralbox}
\small
\begin{multicols}{2}

\textbf{This lesson connects forward to\dots}
\begin{itemize}[itemsep=2pt]
  \item \textbf{Lesson 1.5:} Histograms display \textit{quantitative} data --- bars touch because the scale is continuous.
  \item \textbf{Unit 2:} Segmented bar charts compare two categorical variables.
  \item \textbf{AP Exam:} Free-response often requires you to \textit{construct} a graph \textit{and} write an interpretation.
\end{itemize}

\columnbreak

\textbf{Big Idea --- Write in your own words:}\\
\textit{What is the most important thing to check before trusting any bar chart you see?}\\[2pt]
\writeline\\[1pt]\writeline\\[1pt]\writeline

\end{multicols}
\end{spiralbox}

\vspace{0.1in}

\begin{notesbox}{Additional Notes}
\vspace{0pt}
\writeline\\\writeline\\\writeline\\\writeline\\\writeline\\\writeline
\end{notesbox}

\end{document}
